
\section{Identity and Unity}
\label{sec:identity-and-unity}

\subsection{The Ontological Problem of Unity}

The relation between identity and unity constitutes a further
problem in the ontological analysis of an existent. If identity
concerns the determinate character through which an existent is
the existent that it is, unity concerns a different but closely
related question: in what sense does an existent constitute
a unity?

The two concepts are frequently treated as interchangeable.
An entity is considered identical with itself because it is one,
or its identity is understood as the condition that makes it
a unified entity. Although these formulations point toward
a genuine ontological relation, they risk collapsing two
distinct questions into a single one.

The present ontology maintains that identity and unity are
conceptually distinguishable and must not be treated as
synonymous.

Identity concerns the ontological determination of an existent
as the existent that it is. Unity concerns the condition or
organization through which an existent constitutes a unity.
The latter may involve coherence, integration, continuity
of constitution, or relations among multiple determinations.
However, none of these possibilities should be assumed to
exhaust the meaning of unity before its ontological structure
has been examined.

The distinction is not an assertion that identity and unity
exist independently of one another. They may be inseparably
related in particular existents. The point is that their
relationship cannot be established merely by defining one
through the other.

\subsection{Unity Is Not Identity}

An existent may constitute a unity, but the fact that it is
one does not, by itself, explain what makes it the existent
that it is.

Unity addresses the question of ontological oneness or
coherence. Identity addresses the question of the determinate
character of the existent. These questions may intersect,
but they are not identical.

If unity were sufficient to define identity, every unified
entity would possess an identity in precisely the same
ontological sense, and every difference in identity would
have to correspond to a difference in unity. Neither
conclusion follows from the concept of unity alone.

Two distinct existents may each constitute a unity while
possessing different identities. Conversely, the identity
of an existent cannot be adequately explained merely by
asserting that it constitutes one entity.

Unity therefore does not replace the ontological question
of identity.

The distinction also prevents the assumption that numerical
oneness is sufficient to establish identity. Numerical
identity concerns whether an existent is one and the same
existent across a comparison or a range of conditions.
Unity, by contrast, concerns the condition in which an
existent constitutes a unity. The two questions must be
examined separately.

\subsection{Identity Is Not Unity}

The reverse reduction is equally problematic.

Identity cannot be defined exclusively as the condition
through which an existent is unified. Such a definition
would make the ontological constitution of identity
dependent upon a particular account of unity.

This would create difficulties whenever unity is expressed
through different forms of organization, composition,
or internal plurality.

An existent may possess a determinate identity while
involving multiple components, structures, processes,
or internally differentiated aspects. The presence of
such plurality does not automatically establish that
the existent lacks unity or identity.

Nor does the recognition of unity establish that all
the components or determinations involved share one
and the same identity.

Identity and unity must therefore remain distinct
analytical concepts even when they refer to dimensions
of the same existent.

This distinction is particularly important for an ontology
that does not assume that every existent is simple,
homogeneous, or internally undifferentiated.

\subsection{Unity and Internal Plurality}

Unity must not be equated with the absence of plurality.

An existent may constitute a unity while involving
multiple determinations, differentiated aspects,
components, or relations. Its unity may concern the
organization or coherence of these elements rather
than their numerical or qualitative singularity.

The existence of internal plurality therefore does
not, by itself, contradict unity.

However, the mere presence of multiple elements does
not automatically establish unity either. A collection
of elements may be multiple without constituting a
single ontological unity in the relevant sense.

The problem is consequently to determine what makes
a plurality constitute, participate in, or belong to
a unity.

This question cannot be answered simply by counting
components or identifying their similarities. It
requires examination of the ontological relations
through which the plurality is organized or integrated.

The present section does not assume that every form
of internal plurality is compatible with every form
of unity. Rather, it rejects the general equation
of unity with simplicity and identifies the relation
between plurality and unity as a distinct ontological
problem.

\subsection{Unity and Composition}

One possible basis of unity is composition.

A composed existent may involve multiple constituents
whose relations contribute to the constitution of
a larger entity. In such cases, unity may be associated
with the organization of components into a more
comprehensive whole.

Nevertheless, composition cannot be identified
with unity without qualification.

The presence of components does not necessarily
establish that they form a unified existent. Conversely,
an account of unity should not assume that every
unified existent must be composed of independently
identifiable parts.

It is therefore necessary to distinguish between
the unity of a composed entity and the broader
ontological problem of unity.

The unity of composition concerns a particular
way in which plurality may be organized into a whole.
It does not necessarily provide a universal account
of all forms of unity.

Furthermore, the identity of a composed existent
cannot automatically be reduced to the identities
of its constituents or to the sum of their properties.
The ontological status of the whole and its relation
to its constituents require independent examination.

Whether a composed existent possesses an identity
that is irreducible to the identities of its components
remains a question to be addressed through the broader
analysis of identity and non-reduction.

\subsection{Unity and Structural Coherence}

Unity may also be considered in relation to structural
coherence.

Structural coherence concerns the organization
and interdependence of determinations within an
existent. It differs from a mere aggregation of
elements because the elements may participate in
an organized configuration.

Such organization may contribute to the unity
of an existent, but structural coherence and unity
should not be treated as identical.

A structure may be described through relations
among its components, whereas unity concerns the
ontological status of the existent as a unity.

The distinction is necessary because a description
of structural relations does not, by itself, settle
whether the structure constitutes one existent,
several interrelated existents, or a particular
configuration of entities.

Likewise, an existent's unity may involve aspects
that cannot be exhaustively represented through
a structural description.

The relationship between structural organization
and unity must therefore be investigated without
reducing the latter to the former.

\subsection{Unity of Source and Unity of Constitution}

The concept of unity may refer to different
ontological relations that should not be conflated.

Unity of source concerns a relation of unity
associated with a common origin or source.

Unity of constitution concerns the way in which
an existent is constituted as a unity through
its determinations and their relations.

These formulations identify different questions.
A common source does not automatically establish
that all things originating from it constitute
one and the same existent. Likewise, a unified
constitution does not necessarily establish
a common source.

The distinction prevents an unwarranted inference
from origin to identity or from constitution
to origin.

The ontological significance of a source may
be substantial, but its relation to the identity
and unity of what originates from it must be
established independently.

This section therefore recognizes unity of source
and unity of constitution as distinct analytical
possibilities without assuming that either provides
a complete account of unity.

\subsection{Unity and Individuation}

Unity is also related to individuation, but the
two concepts address different aspects of the
ontological constitution of an existent.

Unity concerns the condition through which
something constitutes a unity. Individuation
concerns the constitution or determination of
something as a distinct existent.

An existent may be unified without the nature
of its individuation being fully explained.
Similarly, the fact that something is individuated
does not, by itself, establish the specific
form of unity it possesses.

The distinction becomes especially important
when considering composed entities and entities
with internal plurality.

An entity may constitute a unity through a
particular organization while its individuation
depends upon conditions that cannot be reduced
to that organization alone.

Conversely, the individuation of an existent
may be established without assuming that it
is simple or internally homogeneous.

Unity and individuation are therefore related
but non-equivalent ontological concepts.

Their precise relation requires an examination
of the conditions under which a plurality
constitutes a single existent and how that
existent is distinguished from others.

\subsection{Unity and Multiplicity}

The distinction between unity and multiplicity
must be carefully maintained.

Multiplicity refers to the presence or condition
of more than one element, entity, determination,
or aspect within a specified ontological context.
Unity refers to the condition through which
something constitutes a unity.

These concepts are not necessarily contradictory
because they may concern different levels or
dimensions of an existent.

An existent may be one in a particular ontological
sense while involving multiplicity in another.
For example, the unity of a composed entity
may coexist with the plurality of its constituents.

However, such coexistence cannot be presumed
in every case. The relevant ontological relations
must be established rather than inferred from
the general compatibility of unity and plurality.

It is also important to distinguish multiplicity
from plurality. These terms may overlap in
ordinary usage, but they need not identify
the same ontological problem.

Multiplicity may concern the presence of
multiple entities or elements, whereas
plurality may refer more broadly to internal
differentiation or the coexistence of distinct
determinations within an existent.

The precise distinction between these concepts
requires further analysis and should not be
settled merely through terminological convention.

\subsection{Unity and the Question of Multiple Identities}

The distinction between unity and identity
raises a further question: can a single existent
possess or involve multiple identities?

The present analysis does not assume that
the unity of an existent necessarily entails
the existence of exactly one identity in every
possible ontological sense.

Nor does it establish that multiple identities
must exist whenever an existent involves
internal plurality or structural differentiation.

The relation between the unity of an existent
and the possible multiplicity of its identities
remains an open ontological problem.

A premature answer would risk confusing several
different questions: whether an existent is
numerically one, whether it constitutes a unity,
whether it contains multiple components, and
whether distinct identity determinations may
be associated with it.

These questions must be distinguished before
any conclusion about the multiplicity of
identity can be justified.

The possibility of multiple identities associated
with one existent is therefore retained as an
open question rather than resolved by the
definition of unity.

\subsection{The Limits of Reducing Unity to Identity}

The analysis establishes several limits on
the relationship between identity and unity.

First, unity cannot be treated as a synonym
for identity. The fact that an existent
constitutes a unity does not exhaust the
question of what makes it the existent
that it is.

Second, identity cannot be defined merely
as the unity of an existent. Such a definition
would make identity dependent upon a
particular interpretation of unity.

Third, unity cannot be equated with simplicity,
homogeneity, or the absence of internal plurality.

Fourth, composition and structural coherence
may contribute to unity without necessarily
providing exhaustive definitions of it.

Fifth, unity of source must be distinguished
from unity of constitution.

Sixth, unity and individuation must remain
analytically distinct, even where they are
ontologically related.

Finally, the unity of an existent does not
settle the open question of whether that
existent may involve multiple identities.

These constraints preserve the conceptual
independence of identity and unity while
allowing their ontological relations to be
investigated.

\subsection{Concluding Formulation}

Identity and unity are distinct but related
dimensions of the ontological constitution
of an existent.

Identity concerns the determinate character
through which an existent is the existent
that it is. Unity concerns the condition
through which an existent constitutes
a unity.

Neither concept can be reduced to the other.
Unity does not, by itself, explain identity,
and identity does not, by itself, provide
a complete account of unity.

Unity may coexist with internal plurality,
composition, and structural differentiation.
Nevertheless, the conditions under which
these forms of plurality constitute a unity
must be examined rather than assumed.

The central principle established in this
section is therefore:

\begin{quote}
An existent's unity and its identity are
ontologically related but conceptually
non-equivalent. The unity of an existent
does not exhaust its identity, and the
identity of an existent does not eliminate
the possibility of internal plurality.
\end{quote}

This principle provides the basis for
examining identity and distinction while
preserving the difference between being
one, being determinate, and being distinct.
